\chapter{El \texttt{SpaceProfile}: cerrar la pila}
\label{ch:spaceprofile}

Un \cod{OdrlProfile} define un vocabulario, pero no dice dónde se aplica. Eso lo dice el \cod{SpaceProfile}: qué
perfil gobierna cada slot, con qué obligación y en qué gate se comprueba. Es además el documento que lleva la
\textbf{versión del espacio}, que es la que publica el hub.

\section{El documento}

\begin{lstlisting}[style=yaml]
dsl: alexandria.profiles/v1
kind: SpaceProfile

space:
  id: https://flexgrid.example/spaces/flexgrid
  version: 1.0.0
  title:
    en: Flexgrid Energy Data Space
    es: Espacio de datos energético Flexgrid
  operator: https://flexgrid.example/party/authority
  base_url: https://flexgrid.example/spaces/flexgrid

lifecycle:
  grace_window: P30D
  propagation_window: P2D
  freshness:
    ttl: PT24H
    max_staleness: P7D
  degraded:
    policy_intake: reject
    dsp_message: continue

bindings:
  - slot: odrl
    profile: https://flexgrid.example/odrl/profile/metering
    source: internal
    version: "^1.0"
    obligation: mandatory
    conformance:
      - { role: json-schema, applies_to: policy_intake }
      - { role: shacl, applies_to: policy_intake }

  - slot: dsp
    profile: https://w3id.org/dspace/2025/1
    source: vendored
    version: "=2025-1"
    obligation: mandatory
    conformance:
      - { role: json-schema, applies_to: dsp_message }
\end{lstlisting}

\section{La sección \texttt{lifecycle}}

Estos parámetros viajan dentro del manifiesto firmado de cada versión, así que todos los clientes aplican los
mismos valores. Lo que no se escribe toma el valor por defecto del hub.

\noindent\begin{tabularx}{\linewidth}{@{}L{5.1cm}C{1.5cm}Y@{}}
\toprule
\textbf{Parámetro} & \textbf{Defecto} & \textbf{Significado} \\
\midrule
\cod{propagation_window} & \cod{P2D} & Tiempo entre que se publica una versión y que empieza a valer. Da margen para sincronizar. \\
\cod{grace_window} & \cod{P30D} & Cuánto convive una versión de ruptura con su predecesora. \\
\cod{min_grace_window} & \cod{P7D} & Gracia mínima que puede pedir un \emph{steward} al aprobar. Es del hub: el \cod{SpaceProfile} no puede cambiarla. \\
\cod{freshness.ttl} & \cod{PT24H} & Validez de un listado. El hub lo reemite al menos cada mitad de este tiempo. \\
\cod{freshness.max_staleness} & \cod{P7D} & Cuánto aguanta un cliente con un listado caducado antes de cerrarse. \\
\cod{sync_interval} & \cod{PT6H} & Intervalo máximo entre sincronizaciones del cliente. \\
\cod{clock_skew} & \cod{PT5S} & Diferencia de reloj tolerada. \\
\cod{lease_margin} & \cod{PT10S} & Margen que el cliente añade alrededor de cada frontera de vigencia. \\
\cod{max_concurrent_versions} & \cod{4} & Máximo de versiones vigentes a la vez. \\
\cod{degraded} & \texttt{\{\}} & Por gate, qué hacer con el listado caducado: \cod{continue} o \cod{reject}. \\
\bottomrule
\end{tabularx}

\medskip
En Flexgrid, \cod{degraded} dice que si el cliente pierde el contacto con el hub sigue validando la envoltura
DSP con lo último que tenía, pero deja de admitir políticas nuevas. El capítulo~\ref{ch:resiliencia} lo provoca.

El compilador comprueba los invariantes entre parámetros. Por ejemplo, con una gracia de tres días y un
\cod{max_staleness} de una hora:

\begin{lstlisting}[style=out]
error   LIFECYCLE_GRACE_TOO_SHORT     s-lifecycle.yaml#2:129 /lifecycle/grace_window grace_window must be at least min_grace_window
error   LIFECYCLE_STALENESS_LT_TTL    s-lifecycle.yaml#2:133 /lifecycle/freshness/max_staleness max_staleness must be greater than ttl
\end{lstlisting}

\section{Los \emph{bindings}}

\noindent\begin{tabularx}{\linewidth}{@{}L{2.6cm}Y@{}}
\toprule
\textbf{Campo} & \textbf{Significado} \\
\midrule
\cod{slot} & \cod{odrl}, \cod{dcat}, \cod{dsp} o \cod{dcp}. Solo un \emph{binding} por slot. \\
\cod{profile} & IRI del perfil que gobierna el slot. \\
\cod{source} & \cod{internal}: un \cod{OdrlProfile} del mismo \emph{bundle}. \cod{vendored}: una copia local de una especificación externa. \\
\cod{version} & Rango SemVer si es interno; versión exacta con \cod{=} si es vendorizado. \\
\cod{obligation} & \cod{mandatory}, \cod{recommended} u \cod{optional}. \\
\cod{via} & Cadena de perfiles intermedios, para perfiles de perfiles. \\
\cod{conformance} & Lista de pares \cod{role} y \cod{applies_to}: qué artefacto se usa en qué gate. \\
\bottomrule
\end{tabularx}

\subsection{Obligación}

La obligación decide qué pasa cuando un mensaje no conforma:

\begin{itemize}[leftmargin=*, itemsep=2pt]
  \item \cod{mandatory}: el gate falla. Un \emph{binding} obligatorio tiene que tener al menos un gate; si no,
        no habría forma de verificarlo.
  \item \cod{recommended}: el gate pasa, y el veredicto lleva el código informativo
        \cod{VERDICT_RECOMMENDED_NONCONFORMANT}.
  \item \cod{optional}: no se evalúa.
\end{itemize}

\subsection{Rangos de versión}

Un \emph{binding} interno no fija una versión del perfil: da un rango, y el compilador elige la mayor versión
del \emph{bundle} que lo cumple.

\noindent\begin{tabularx}{\linewidth}{@{}L{2.6cm}Y L{4.6cm}@{}}
\toprule
\textbf{Rango} & \textbf{Acepta} & \textbf{Ejemplo} \\
\midrule
\cod{=1.2.3} & Exactamente esa versión. & \\
\texttt{\textasciicircum{}1.0} & Misma versión mayor, igual o superior. & \cod{1.0.0}, \cod{1.4.2}; no \cod{2.0.0} \\
\texttt{\textasciicircum{}0.3} & Con mayor 0, misma versión menor. & \cod{0.3.0}, \cod{0.3.9}; no \cod{0.4.0} \\
\texttt{\textasciitilde{}1.2} & Misma mayor y menor. & \cod{1.2.0}, \cod{1.2.7}; no \cod{1.3.0} \\
\cod{1} & Sin operador, las partes que faltan son comodines. & cualquier \cod{1.x.y} \\
\bottomrule
\end{tabularx}

\medskip
Las versiones preliminares, como \cod{1.1.0-rc1}, no satisfacen ningún rango.

\subsection{Roles y gates}

\cod{conformance} une un \textbf{rol} de artefacto con un \textbf{gate}. Un perfil interno ofrece los roles
\cod{json-schema} y \cod{shacl}. Un perfil vendorizado ofrece los que declare su \cod{vendor.json}.

\noindent\begin{tabularx}{\linewidth}{@{}L{5.2cm}Y C{2.4cm}@{}}
\toprule
\textbf{Gate} & \textbf{Qué se valida} & \textbf{¿Lo aplica el cliente?} \\
\midrule
\cod{dsp_message} & La envoltura del mensaje DSP. & \si \\
\cod{policy_intake} & La política de una oferta, una petición o un acuerdo. & \si \\
\cod{catalog_entry} & Las respuestas de catálogo. & \si \\
\cod{pdp_load} & La política al cargarse en el motor de decisión. & \no \\
\cod{credential_presentation} & Las credenciales presentadas. & \no \\
\bottomrule
\end{tabularx}

\begin{trampaHabitual}[Dos gates del lenguaje no los aplica nadie]
El lenguaje admite \cod{pdp_load} y \cod{credential_presentation}, y un \emph{binding} en ellos compila. Pero
\cod{alexandria-client} solo conoce \cod{dsp_message}, \cod{policy_intake}, \cod{catalog_entry} y
\cod{agreement_admission}. Un gate que ningún componente evalúa es documentación, no control.
\end{trampaHabitual}

\begin{trampaHabitual}[Un gate obligatorio cubierto solo por SHACL]
Hoy ningún gate ejecuta SHACL: el cliente solo ejecuta \cod{json-schema}. Si un \emph{binding} \cod{mandatory}
enlaza \emph{únicamente} \cod{shacl} a un gate, ese gate acepta cualquier mensaje. El compilador lo avisa:

\begin{lstlisting}[style=out]
warning GATE_SHACL_ONLY   shacl-only.yaml#2:145 /bindings/0/conformance gate policy_intake of mandatory slot odrl is covered only by shacl, which no gate executes yet: it will accept any message
\end{lstlisting}

Por eso Flexgrid enlaza los dos roles a \cod{policy_intake}: el esquema valida hoy, y las \emph{shapes} quedan
publicadas para cuando haya motor.
\end{trampaHabitual}

\section{Perfiles vendorizados}

DSP no es un perfil que escriba Flexgrid: es una especificación externa que el espacio \textbf{fija}. Como nada se
resuelve contra la red, hay que tener una copia local. El compilador la busca en
\texttt{<vendor>/\allowbreak{}<slot>/\allowbreak{}<versión>/\allowbreak{}vendor.json}:

\begin{lstlisting}[style=json]
{
  "profile": "https://w3id.org/dspace/2025/1",
  "version": "2025-1",
  "artifacts": [
    { "path": "dsp-messages.schema.json", "role": "json-schema",
      "media_type": "application/schema+json" },
    { "path": "context.jsonld", "role": "context",
      "media_type": "application/ld+json",
      "url": "https://w3id.org/dspace/2025/1/context.jsonld" }
  ]
}
\end{lstlisting}

El campo \cod{profile} tiene que coincidir con el del \emph{binding}. Cada artefacto declara su ruta, su rol y
su tipo. El campo \cod{url} es importante en los contextos JSON-LD: es la URL con la que los mensajes los citan,
y el cliente solo resolverá las URL que el paquete declare.

\begin{notaPractica}[La copia de este tutorial es reducida]
Los dos ficheros de \cod{examples/vendor/dsp/2025-1/} están escritos a mano para el laboratorio: un esquema que
comprueba los campos obligatorios de tres mensajes y un contexto con los términos que usan los ejemplos. En un
espacio real se copian ahí los ficheros oficiales de DSP y se fija su versión.
\end{notaPractica}

Sin la opción \cod{-vendor}, el \emph{binding} no resuelve:

\begin{lstlisting}[style=sh]
../../../bin/alexandriac validate dsl/flexgrid-1.0.0.yaml
\end{lstlisting}
\begin{lstlisting}[style=out]
error   BINDING_UNRESOLVED   flexgrid-1.0.0.yaml#2:148 /bindings/1 no vendored copy of https://w3id.org/dspace/2025/1 2025-1 for slot dsp
alexandriac: 1 errors
\end{lstlisting}

\section{El \emph{bundle} y su compilación}

Un \textbf{bundle} es un flujo YAML: un \cod{SpaceProfile} y los \cod{OdrlProfile} que enlaza, separados por
\cod{---}. \cod{examples/dsl/flexgrid-1.0.0.yaml} es exactamente el perfil del capítulo anterior seguido del
espacio de este.

\begin{lstlisting}[style=sh]
../../../bin/alexandriac compile -o /tmp/flexgrid -vendor vendor dsl/flexgrid-1.0.0.yaml | sort
\end{lstlisting}
\begin{lstlisting}[style=out]
/tmp/flexgrid/profiles/dsp/vendored/2025-1/context.jsonld
/tmp/flexgrid/profiles/dsp/vendored/2025-1/dsp-messages.schema.json
/tmp/flexgrid/profiles/odrl/metering/1.0.0/context.jsonld
/tmp/flexgrid/profiles/odrl/metering/1.0.0/prof.ttl
/tmp/flexgrid/profiles/odrl/metering/1.0.0/schema.json
/tmp/flexgrid/profiles/odrl/metering/1.0.0/shapes.ttl
/tmp/flexgrid/profiles/odrl/metering/1.0.0/spec.html
/tmp/flexgrid/profiles/odrl/metering/1.0.0/vocab.ttl
/tmp/flexgrid/space/prof.ttl
/tmp/flexgrid/space/shacl-union.ttl
/tmp/flexgrid/space/space-lock.json
/tmp/flexgrid/space/spec.html
\end{lstlisting}

Aparecen dos cosas nuevas: los ficheros vendorizados, copiados bajo \texttt{profiles/\allowbreak{}<slot>/\allowbreak{}vendored/\allowbreak{}<versión>/\allowbreak{}}, y
cuatro artefactos del espacio.

\noindent\begin{tabularx}{\linewidth}{@{}L{4.8cm}Y@{}}
\toprule
\textbf{Artefacto del espacio} & \textbf{Contenido} \\
\midrule
\cod{space/space-lock.json} & La resolución: versión elegida de cada \emph{binding}, sus artefactos con SHA-256 y gates, los parámetros de ciclo de vida y los perfiles declarables. \\
\cod{space/prof.ttl} & Manifiesto PROF del espacio: un \cod{prof:isProfileOf} por \emph{binding} y un descriptor por artefacto. \\
\cod{space/shacl-union.ttl} & Las \emph{shapes} de los \emph{bindings} obligatorios, concatenadas por gate. \\
\cod{space/spec.html} & Tabla legible de \emph{bindings}. \\
\bottomrule
\end{tabularx}

\medskip
El \cod{space-lock.json} es el equivalente a un fichero de bloqueo de dependencias. Donde el YAML decía
\texttt{"\textasciicircum{}1.0"}, el bloqueo dice qué se eligió y con qué bytes:

\begin{lstlisting}[style=json]
{ "slot": "odrl", "source": "internal", "obligation": "mandatory",
  "profile": "https://flexgrid.example/odrl/profile/metering",
  "version": "1.0.0",
  "artifacts": [
    { "path": "profiles/odrl/metering/1.0.0/schema.json", "role": "json-schema",
      "media_type": "application/schema+json", "applies_to": ["policy_intake"],
      "sha256": "fa70ada7f8c3e45426e804c1883da0a338d3e718425211ad86626d702cdc5e09" },
    { "path": "profiles/odrl/metering/1.0.0/shapes.ttl", "role": "shacl",
      "media_type": "text/turtle", "applies_to": ["policy_intake"],
      "sha256": "43d5fb036d58eea976c45b4af11575ec0e012162ce2b29ad35185900eb4929c8" },
    … ] }
\end{lstlisting}

\begin{ideaClave}
Cada artefacto lleva los gates de su \emph{binding}. \cod{schema.json} del slot \cod{odrl} se aplica en
\cod{policy_intake} porque el \cod{SpaceProfile} lo dice, no porque el artefacto lo sepa. Cambiar dónde se valida
algo es editar una línea del espacio, sin tocar el perfil.
\end{ideaClave}

\section{Errores propios del espacio}

\noindent\begin{tabularx}{\linewidth}{@{}L{4.6cm}Y@{}}
\toprule
\textbf{Código} & \textbf{Causa} \\
\midrule
\cod{SLOT_DUP} & Dos \emph{bindings} para el mismo slot. \\
\cod{BINDING_UNRESOLVED} & Ninguna versión del \emph{bundle} cumple el rango, o falta la copia vendorizada. \\
\cod{BINDING_BASE_MISMATCH} & Un slot distinto de \cod{odrl} declarado como \cod{internal}: el lenguaje solo tiene documentos \cod{OdrlProfile}. \\
\cod{ARTIFACT_MISSING} & Se pide un rol que el perfil no produce, como \cod{cedar-mapping}. \\
\cod{MANDATORY_UNVERIFIABLE} & Un \emph{binding} obligatorio sin gates. El esquema lo detiene antes, como error \cod{SCHEMA}. \\
\cod{VIA_CYCLE} & \cod{via} contiene el propio perfil. \\
\bottomrule
\end{tabularx}

\medskip
Las salidas reales de los dos primeros:

\begin{lstlisting}[style=out]
error   SLOT_DUP             s-slot-dup.yaml#2:156 /bindings/2/slot slot odrl has more than one binding
error   BINDING_UNRESOLVED   s-unresolved.yaml#2:139 /bindings/0 no release of https://flexgrid.example/odrl/profile/metering satisfies ^2.0 in the bundle
\end{lstlisting}

En todos los mensajes, \texttt{\#2} es el número de documento dentro del flujo YAML, y lo que sigue a los dos puntos
es la línea del fichero.

\begin{retos}
\begin{enumerate}[leftmargin=*, itemsep=2pt]
  \item Cambia la obligación del slot \cod{odrl} a \cod{recommended} y compara el \cod{space-lock.json}.
  \item Añade un \emph{binding} \cod{dcat} vendorizado sin crear su directorio y lee el error. Después crea
        \cod{vendor/dcat/3.0.0/vendor.json} con un artefacto de rol \cod{shacl} y enlázalo a \cod{catalog_entry}.
  \item Sube \cod{max_concurrent_versions} a 1 y piensa qué implicaría para una versión de ruptura.
\end{enumerate}
\end{retos}
